\documentclass[11pt,a4paper]{article}
\usepackage[T1]{fontenc}\usepackage{lmodern,amsmath,amssymb,amsthm,mathtools,booktabs,microtype}
\usepackage[margin=26mm]{geometry}\usepackage[colorlinks=true,linkcolor=blue,urlcolor=blue,citecolor=blue]{hyperref}
\newtheorem{theorem}{Theorem}[section]\newtheorem{lemma}[theorem]{Lemma}\newtheorem{proposition}[theorem]{Proposition}\newtheorem{corollary}[theorem]{Corollary}
\newcommand{\Q}{\mathbb Q}\newcommand{\Z}{\mathbb Z}\DeclareMathOperator{\Tr}{Tr}
\setlength{\emergencystretch}{2em}
\title{A shifted Lucas equation and a quartic in 200 variables}
\author{}\date{15 November 2026; corrected scalar normalization}
\begin{document}\maketitle
\begin{abstract}
A shifted half-integral Lucas label is automatically nonzero, while its
associated norm equation retains the growth needed for the height argument.
This permits replacing the last separately stored Lucas root point in the
204-variable source by a recovered elliptic difference. The resulting
construction has 200 total rational variables and 186 quadratic equations.
Their primitive integral sum of squares has exact degree four and content
one. The argument retains the inherited arithmetic premises of the source;
it does not supply an independent audit of those premises. It establishes
an existence family with an ineffective sufficient coefficient threshold,
not a numerically certified defining member. The two comparison poles remain
independent: this reduction does not require or establish matching pole
square classes.
\end{abstract}

\section{Statement, source, and arithmetic scope}
Every variable count includes the input $X$. For a reduced rational $u/v$, with $v>0$, write $h(u/v)=\log\max(|u|,v)$. For $a+b\rho\in K$ write $\operatorname{Re}(a+b\rho)=a$, and let $\sigma$ denote the nontrivial automorphism of $K$. Put
\[
 K=\Q(\rho),\quad \rho^2=2,\quad d=3-\rho,\quad D=d^2=11-6\rho,
 \quad E:y^2=x^3-Dx,
\]
\[
 R_0=(1+\rho,2\rho),\qquad P=2R_0.
\]
We use the source group assertion $2E(K)=\Z P$, with $P$ non-torsion.
For every $Q\in E(K)$ define $\nu(Q)\in\tfrac12\Z$ by
$2Q=[2\nu(Q)]P$. The kernel of this homomorphism is $E[2](K)$.

The parent is the unbalanced 204-variable system in
\begin{center}\small\nolinkurl{computation/quartic/reciprocal204_v27/unbalanced-quadratic.json}.
\end{center}
The complete source proof is \cite{Parent}. The transformation below fixes
this table, rather than an unspecified presentation with the same count.
The arithmetic inputs retained from it are the group assertion, scalar
height and uniform scalar completeness statements, the full-family
contact-height bound, the pole obstruction and prepared pole realization,
and the fixed-data prime-pattern and compatible real-progression statements.
We also use the same projective-family height theorem and pointed elliptic
independence as the source weighted-trace proof \cite{Weighted}.
Section~\ref{sec:weighted} verifies the changed application of that theorem.
No new $abc$ or Szpiro hypothesis is imposed.

\begin{theorem}\label{thm:main}
Under these inherited arithmetic premises there is a computable family
\[
 P_N\in\Z[X,Y_1,\ldots,Y_{199}],\qquad N\ge1,
\]
with total degree exactly four and integer content one, given as the sum
of 186 primitive integral quadratic squares. There is one integer $N_0$
such that, for every $N\ge N_0$,
\[
 \{x\in\Q:\exists\boldsymbol y\in\Q^{199},\ P_N(x,\boldsymbol y)=0\}=\Z.
\]
The sufficient threshold $N_0$ is ineffective. It is not asserted to equal
the parent's threshold. The construction does not change the separately
certified parameter of the explicit 260-variable family.
\end{theorem}

The balanced family has exactly 3221 nonzero coordinate monomials. Its
coefficient height is determined in Section~\ref{sec:balance}.
The polynomial is specified by the exact row transformation in
Section~\ref{sec:spec}. All fixed powers of 149 are coefficients, not
quantified variables. Computational reconstruction and the mathematical
arithmetic argument are separate verification tasks.

\section{Correction of the unpublished coordinate formula}
The 14 November 2026 package shifted the stored second-ratio slots by
2 instead of $2R$. Its coefficient checker repeated that substitution
and therefore did not verify its physical interpretation. For a physical
ratio $\ell$, the old slot is $R\ell$. The correct substitution is
$L_{2,\mathrm{old}}=(L_2-2R)/2$, giving $L_2/R=2\ell+2$.
The former substitution instead gives $L_2/R=2\ell+2/R$ and does not
supply the integer Lucas label used in its proof. That earlier table is
superseded; no equivalence to it is claimed. The corrected substitution
is affine in retained coordinates, preserves quadratic degree, and is
invertible since the unchanged scalar source forces $R>100$.

\section{A Lucas equation with an automatically nonzero label}
For an integer $A$ define
\[
 U_0(A)=0,\quad U_1(A)=1,\quad U_{j+1}(A)=A U_j(A)-U_{j-1}(A).
\]

\begin{lemma}[Integral norm-one solutions]\label{lem:units}
Let $C\ge3$ be an integer. If integers $t,w$, with $w\ne0$, satisfy
$t^2-Ctw+w^2=1$, then for some $j\ge1$
\[
 |w|=U_j(C),\qquad |t|\in\{U_{j-1}(C),U_{j+1}(C)\}.
\]
Moreover $U_j(C)\ge(C-1)^{j-1}$ for $j\ge1$.
\end{lemma}
\begin{proof}
Put $\epsilon=(C+\sqrt{C^2-4})/2$. The element $t-w\epsilon$
is a norm-one unit of the order $\Z[\epsilon]$. Every positive norm-one
unit in its first real embedding is an integral power of $\epsilon$.
Indeed $\epsilon$ is itself such a unit. There is no unit $\alpha$ in
$1<\alpha<\epsilon$: writing $\alpha=u+v\epsilon$ gives
\[
 0<v=\frac{\alpha-\alpha^{-1}}{\epsilon-\epsilon^{-1}}<1,
\]
contrary to integrality. Dividing any positive unit by a suitable integral
power of $\epsilon$ reduces it to this interval or to 1. Negative units
are obtained by multiplication by $-1$. The identity
$\epsilon^j=U_j(C)\epsilon-U_{j-1}(C)$ and conjugation now give the two
stated possibilities. The excluded case $w=0$ is exactly the zeroth power.
Finally $U_2=C\ge C-1$, and the recurrence with
$U_j\ge U_{j-1}>0$ gives $U_{j+1}\ge(C-1)U_j$.
\end{proof}

\begin{theorem}[Shifted labels]\label{thm:lucas}
Fix an integer $r\ge2$. Let $a,m,t\in\Z$, $c=|a|\ge4$, and put
\[
 A=ra+2,\qquad b=m+1/r.
\]
If
\begin{equation}\label{eq:lucas}
 t^2-A(A-2)tb+(A-2)^2b^2=1,
\end{equation}
then $b\ne0$ and
\begin{equation}\label{eq:growth}
 \log\max(1,|t|)\ge(c-2)\log(rc-3).
\end{equation}
Conversely, for every positive integer $a\ge2$ one has an integer solution
by setting
\begin{equation}\label{eq:lucas-complete}
 A=ra+2,\qquad t=U_{a+1}(A),\qquad
 m=\frac{U_a(A)/a-1}{r}.
\end{equation}
\end{theorem}
\begin{proof}
The number $b=m+1/r$ cannot vanish for integral $m$ and $r\ge2$.
Put $w=(A-2)b=a(rm+1)\in\Z\setminus\{0\}$.
Equation~\eqref{eq:lucas} becomes $t^2-Atw+w^2=1$.
Replacing $w$ by $\operatorname{sgn}(A)w$ reduces to $C=|A|\ge rc-2\ge6$.
Lemma~\ref{lem:units} gives $|w|=U_j(C)$ for $j\ge1$.
Since $C\equiv2$ or $-2\pmod c$, the recurrence yields
\[
 U_j(C)\equiv j\quad\hbox{or}\quad(-1)^{j-1}j\pmod c.
\]
But $c\mid w$, so $c\mid j$ and $j\ge c$.
The smaller possible absolute value of $t$ is $U_{j-1}(C)$, which is
at least $(C-1)^{j-2}\ge(rc-3)^{c-2}$. This proves~\eqref{eq:growth}.

For the converse, $A\equiv2\pmod{ra}$ gives $U_a(A)\equiv a\pmod{ra}$;
hence the displayed $m$ is integral and $(A-2)b=U_a(A)$.
The identity
\[
 U_{j+1}(A)^2-AU_{j+1}(A)U_j(A)+U_j(A)^2=1
\]
follows by induction from the recurrence, proving~\eqref{eq:lucas}.
For the application $r=2$, $a\ge3$ implies $m>0$.
\end{proof}

We use $r=2$. The important change is $A=2a+2$, not $A=2a$:
\[
                 b=m+\tfrac12,\qquad (A-2)b=a(2m+1).
\]
It preserves integrality of the entire residual and removes the need for a
separate point certificate imposing a nonzero fourth integer label.

\section{The unchanged scalar source}
Write $R=x_{118}$ and $R_2=x_{117}$ in the parent's coordinates.
The retained scalar equations imply
\begin{equation}\label{eq:scalar}
 \begin{gathered}
 \beta=51X^2+2504,\qquad
 r_0=\beta(\beta-1)(\beta-2401)(32\beta+343),\\
 (1+149^2r_0)R=149^{2N+1}r_0,\qquad R_2=R^2.
 \end{gathered}
\end{equation}
Since $\beta\ge2504$, one has $r_0\ge1$ and $R>100$.
The scalar point is still $Q_2=[a]P$ with nonzero integer $a$.
The scalar proof uses that physical point and the reference point, not the
new value assigned to its comparison label. Precisely its retained conclusion is
\begin{equation}\label{eq:scalar-bound}
 c=|a|\ge149^{L_N},\qquad
 h(\beta)\le C_3\,149^{-L_N}c(1+\log c/3),
 \quad L_N=\left\lfloor\frac{9(2N-1)}4\right\rfloor+1,
 \quad C_3=2^{330770}.
\end{equation}
Uniform scalar completeness allows positive $a$ as large as required,
after its input-dependent finite preparation, for every $N$ above one
fixed scalar threshold. We use exactly the order of choices in the parent:
choose this scalar point before the Lucas points and the comparison primes.

The reference equations recover $H,H-P$ and the nonzero quantities
\[
 S=y(H)/x(H),\qquad T_H=y(H-P)/x(H-P).
\]
Its index avoids 0 and 1. These assertions and the scalar nonvanishing
arguments remain those of the parent. In particular the index $a$ in
\eqref{eq:scalar-bound} is not replaced by $A$.

\section{The recovered fourth difference and its ratios}
For $Q\in E\setminus\{(0,0)\}$ write
\[
 f(Q)=1/x(Q),\qquad h(Q)=-y(Q)/x(Q)^2,
 \quad f(O)=h(O)=0.
\]
The reciprocal curve is $h^2=f-Df^3$. It is an ordinary affine model
of the original curve translated by $(0,0)$.
Keep the already stored squares $F,G$ and introduce five field coordinates
$f,h,g,k,s$. Require
\begin{equation}\label{eq:pair}
 \begin{split}
 F=f^2,\qquad G=g^2,\qquad h^2=f-DfF,\\
 s(f-g)=h+k,\qquad s(h-k)=1-D(F+fg+G).
 \end{split}
\end{equation}
The fourth root point is recovered, without new output coordinates, by
\begin{equation}\label{eq:recover}
 x_4=s^2+D(f+g),\qquad y_4=Dh-s(Df+x_4).
\end{equation}

\begin{lemma}\label{lem:pair}
Every rational solution of~\eqref{eq:pair} gives two points
$U_4,V_4\in E(K)\setminus\{(0,0)\}$ with $U_4\ne V_4$, and
$Q_4=U_4-V_4=(x_4,y_4)$. Conversely every such pair gives the unique
secant or tangent parameter $s$. Identity numerator points and finite
nonzero two-torsion differences are included.
\end{lemma}
\begin{proof}
Put $r_1=s(f-g)-h-k$, $r_2=s(h-k)-1+D(f^2+fg+g^2)$,
$r_3=h^2-f+Df^3$. Direct polynomial identities are
\[
 k^2-g+Dg^3=r_3-(f-g)r_2+(h-k)r_1,
\]
\[
 y_4^2-x_4^3+Dx_4=D^2r_3-D(Df+x_4)(r_2-sr_1).
\]
For $f\ne g$, the chord law on the translated model gives precisely
\eqref{eq:recover}. If $f=g$, the first linear slope equation gives
$k=-h$. The possibility $h=0$ would imply $f=0$ or $Df^2=1$, while
the other slope equation would imply $1-3Df^2=0$, a contradiction.
Hence $h\ne0$ and $2sh=1-3Df^2$, the tangent formula. These cases
also prove uniqueness and the converse. In particular no division by
$f$, $g$ or a point ordinate was used to exclude an identity component.
\end{proof}

Keep the two old fourth ratio slots $L_4,L_4^0$ and put
$M=L_4-\tfrac12$, $M^0=L_4^0-\tfrac12$. Add four field quadrics
\begin{equation}\label{eq:ratios}
 \begin{split}
 hM+2fS=0,&\qquad (1-DF)M+2hS=0,\\
 kM^0+2gT_H=0,&\qquad (1-DG)M^0+2kT_H=0.
 \end{split}
\end{equation}
On the nonidentity domain they express
\[
 L_4=2z(U_4)/z(H)+\tfrac12,\qquad
 L_4'=2\sigma(z(V_4)/z(H-P))+\tfrac12,\qquad z=-x/y.
\]
At an identity numerator the second row forces the corresponding ratio
to $1/2$. A nonidentity two-torsion numerator would have $h=0,f\ne0$
and force $S=0$, or the analogous contradiction at $T_H$; it supplies
no spurious component.

In the retained second root block make the invertible substitutions
\begin{equation}\label{eq:affine}
 L_{2,\rm old}=(L_2-2R)/2,\qquad
 L_{2,\rm old}^0=(L_2^0-2R)/2.
\end{equation}
The physical point $Q_2$ and all its scalar equations are unchanged.
The slots $L_2,L_2^0$ store $R$ times their physical comparison
ratios. Their physical comparison label is now $A=2a+2$; this is not
an assertion that the stored slot itself equals that label modulo a prime. The comparison-factor rows outside
that root block continue to use their displayed slots $L_2,L_2^0$.
The added multiples of $R$ cancel in cross differences. The stored
cross-square row divides their squared difference by $R^2$, and the
Lucas product rows divide the products involving $L_2$ by $R$. Thus
both use the physical ratio $L_2/R$, not the stored slot $L_2$. The extracted labels are
\[
 n=\nu(Q_1)\in\tfrac12\Z,\quad A=2a+2,\quad
 t=2\nu(Q_3)\in\Z,\quad b=2\nu(Q_4)+\tfrac12=m+\tfrac12,
 \qquad m\in\Z.
\]

\section{The changed trace-height application}\label{sec:weighted}
Replace the parent's aggregate by
\begin{equation}\label{eq:aggregate}
 t_1=\lambda+T,\qquad
 T=\tfrac14\Tr\bigl(x(Q_1)+x(Q_3)+x(Q_4)+R^2x(Q_2)\bigr).
\end{equation}
It is still represented by one rational quadratic row: the new $x(Q_4)$
is quadratic by~\eqref{eq:recover}, and the second root contribution
$R^2x(Q_2)$ uses the same stored difference of square-root coordinates
as in the parent.

\begin{lemma}\label{lem:height}
Put $B_0=\max(1,|n|,|a|,|t|,|m|)$. There are fixed constants
$A_0>0,J_0\ge0$, independent of $N$, the input, and the witnesses, such that
\[
 T>0,\qquad B_0^2\le A_0\bigl(h(R)+h(T)\bigr)+J_0.
\]
\end{lemma}
\begin{proof}
At either real embedding every finite point has $x\ge-d$ in that
embedding. The nonzero doubled scalar point has $x>d$. Thus
$4T>6R^2-18>0$.

For the height assertion apply the same projective-family theorem as
\cite{Weighted}, to the graph of the sum in~\eqref{eq:aggregate}, with
parameters $(R,T)$. The only change in its weights is the partition
$(1,1,R^2,R^2)$ into $(1,1,1,R^2)$. We check the required
nonconstancy on abelian translates rather than infer it from positivity
alone. First fix all index-zero coordinates: they belong to the finite
set of two-torsion points. This gives finitely many strata. In a stratum,
the remaining pointed coordinates are nonzero rational multiples of the
fixed pointed generator in each of the two conjugate blocks, up to torsion.

Along an elliptic subgroup direction, a common pole of two coordinates
in one block implies that the quotient of their nonzero endomorphism
coefficients is the quotient of their nonzero rational indices. Indeed
eliminating the subgroup parameter yields an endomorphism of a non-torsion
point that is torsion; that endomorphism must be zero. The principal
coefficients of their abscissa poles therefore have a common nonzero
factor times positive rational squares. Multiplying by the positive
weights $1$ or $R^2$ cannot cancel their sum. Common poles between the
conjugate blocks are excluded by the same independence of the two pointed
images over $\Z[i]$ proved and used in \cite{Weighted}. Directions in
index-zero coordinates have already been removed by fixing that stratum.
Every positive-dimensional abelian subvariety of the relevant power of the fixed CM elliptic curve contains an elliptic subgroup direction. Hence no positive-dimensional abelian translate through an actual tuple
can make this weighted sum constant. Boundary components of the projective
graph do not affect this argument, since a translate through a tuple in
the function domain meets that domain densely.

The ambient product, its ample embedding, and the fixed finite-rank group
are unchanged. The height of the parameterized coefficients $[1:R^2]$
is $2h(R)$, and the target parameter contributes $h(T)$, up to a fixed
constant. The family theorem therefore bounds the product canonical
height by a constant times $h(R)+h(T)+1$. The canonical heights of the
non-torsion coordinates are fixed positive multiples of their squared
indices; the factors two in $t,m$ change only the constants. Taking
maximum constants over the finitely many torsion strata proves the claim.
This is the same height-theorem input as in the parent, with its ineffective
additive constant retained.
\end{proof}

\section{Local extraction and the two integer residuals}
Keep every contact, detector-factor, trace, and pole row. The contacts are
\[
 \Theta_0=\{0,-1/2,-1,1/2399,-32/407\},\qquad
 \Theta=\Theta_0\cup(2-\Theta_0),
\]
and the exceptional primes are $\Sigma=\{p:p\le76176\}\cup\{2029147\}$.
Let $M_j$ be the parent's odd contact radical at $t_j$, for
$t_0=\lambda,t_1=\lambda+T$. Its retained full-family bound is
\[
                         h(t_j)\le H_2M_j^{7/6}.
\]
In the parent notation the five local factors are
\begin{align*}
 1+q^5C_2^2+q^2D_H^2+q^4B_1^4,&\quad
 1+q^5G^2+q^2A_3^2+q^4B_0^4,\\
 1+q^5I^4+q^3A_2^2+q^2A_4^2+q^4B_3^4,&\\
 1+q^5C_3^2+q^2C_1^4+q^4B_4^4,&\quad
 1+q^5C_4^2+q^2A_1^2+q^4B_2^4.
\end{align*}
Here $G$ is the displayed Lucas residual in the ratio slots, not the
reciprocal square $G$ of Section 4; these are separate named source slots.
The tests $A_i,B_i$ are reciprocal abscissas, $C_i$ are cross ratios,
and $I$ is the input comparison. Their valuation classes and the pole
obstruction are unchanged. Zero tests are allowed. In particular the
recovered reciprocal abscissas still have even valuation, including the
identity convention.

The parent's integral-logarithm argument applies to the unshifted ratios
obtained from~\eqref{eq:ratios} and~\eqref{eq:affine}. All consumed primes
are odd, so the factors two and the shift $1/2$ are units or integral.
At least one reference index $h,h-1$ is a $p$-adic unit. One exact
logarithmic ratio is therefore integral; positive cross precision transfers
integrality to the other. Only then is relative precision four used as
absolute precision. Consequently the extracted labels are precisely
$n,A,t,b$ above, with cross and Lucas precision three and input precision
one. Thus
\begin{equation}\label{eq:divisibility}
 M_j^3\mid\Delta\quad\hbox{if }\Delta\ne0,\qquad
 M_j\mid J\quad\hbox{if }\beta\ne n,
\end{equation}
where
\begin{equation}\label{eq:delta}
 \Delta=t^2-(2a+2)a(2m+1)t+a^2(2m+1)^2-1\in\Z
\end{equation}
and $J$ is the positive reduced numerator of $\beta-n$.
These are separate divisibilities for each detector; coprimality between
$M_0$ and $M_1$ is not asserted. The elementary bounds are
\begin{equation}\label{eq:bounds}
 |
 \Delta|\le24B_0^4,\qquad J\le4B_0e^{h(\beta)}.
\end{equation}
For the second bound multiply a reduced expression for $\beta$ by two
and use $2n\in\Z$; taking the reduced numerator can only decrease it.

\section{Uniform soundness}
The height of the rational function in~\eqref{eq:scalar} satisfies
\[
 h(R)\le4h(\beta)+(2N+1)\log149+C.
\]
The degree four in $\beta$ can be checked using homogeneous numerator
and denominator of the same degree; the coefficient contribution from
$1+149^2r_0$ is fixed apart from the displayed power of 149.
Also $N\le C'(1+\log c)$ by~\eqref{eq:scalar-bound}.

Suppose first that $\Delta\ne0$. Equations~\eqref{eq:divisibility},
\eqref{eq:bounds}, and the contact-height bound give
$h(T)\le C|\Delta|^{7/18}+C$ because
$h(T)\le h(t_0)+h(t_1)+\log2$.
Lemma~\ref{lem:height} and~\eqref{eq:scalar-bound}, with $c\le B_0$, yield
\[
 B_0^2\le C B_0^{14/9}+C B_0(1+\log B_0)+C.
\]
All constants here are independent of $N$. The right-hand side is
$o(B_0^2)$, so $B_0$ is bounded. This contradicts
$c\ge149^{L_N}$ once $N$ exceeds a fixed threshold. Hence $\Delta=0$.
Theorem~\ref{thm:lucas}, with $r=2$, now gives
\[
                 \log B_0\ge(c-2)\log(2c-3).
\]
No nonzero guard on $m$ or $Q_4$ was needed for this conclusion:
$b=m+1/2$ already made the norm coefficient nonzero.

If $\beta\ne n$, the same argument at order one gives
\[
 B_0^2\le C(B_0e^{h(\beta)})^{7/6}+C h(R)+C.
\]
The last two terms, bounded uniformly by $C B_0(1+\log B_0)+C$, can
be absorbed for large $B_0$, which is ensured by large $N$.
Taking logarithms then gives
\[
 \log B_0\le C+\frac75h(\beta)
 \le C+\frac75C_3\,149^{-L_N}c(1+\log c/3).
\]
Choose $N$ sufficiently large that $(7/5)C_3 149^{-L_N}\le1$.
The upper bound $C+c+(c/3)\log c$ contradicts the lower bound above
for all sufficiently large $c$. This proves $\beta=n\in\tfrac12\Z$.
Writing $X=u/v$ in lowest terms with $v>0$, the integrality of
$2\beta=102X^2+5008$ forces $v^2\mid102$. Since 102 is squarefree,
$v=1$. All large-$N$ choices made here are fixed independently of the
input and the witnesses. They remain ineffective because the height
constant in Lemma~\ref{lem:height} remains ineffective.

\section{Completeness and the order of choices}
Fix an integer $X$ and take $N$ above the parent's uniform scalar threshold.
Use the unchanged scalar preparation to choose a positive integer $a\ge3$
and $Q_2=[a]P$. Set
\[
 A=2a+2,\qquad t=U_{a+1}(A),\qquad
 m=(U_a(A)/a-1)/2>0,
\]
and choose
\[
 Q_1=[\beta]P,\qquad Q_3=[t]R_0,\qquad Q_4=[m]R_0.
\]
Theorem~\ref{thm:lucas} gives the required Lucas identity exactly.
Fix these points and the positive rational scale $T$ before selecting
comparison primes. The scalar data and its continuity conditions remain
fixed at this stage.

The parent completeness proof already uses the half-generator $R_0$ for
$Q_3$. Its half-generator congruence lemma says that the prescribed
Gaussian congruences make the same primary and opposite annihilating
order moduli kill $R_0$. Using an additional integer multiple $[m]R_0$
therefore introduces no additional congruence modulus or prime-pattern
hypothesis. Apply the same fixed-data prime selection and compatible CRT
to obtain the reference index $h$. Put $Z=[h]R_0$; then
\[
 H=2Z=[h]P,\quad
 U_1=[2\beta]Z,\quad U_2=[2a]Z,\quad U_3=[t]Z,\quad U_4=[m]Z,
 \qquad V_i=U_i-Q_i.
\]
At the selected primary place $Z$ is formal, and at the opposite place
$Z-R_0$ is formal. Hence all the required $U_i,V_i$ have the same
formal depths as in the parent. The exact logarithmic ratios, after
the factors and shifts above, are $\beta,A,t,m+1/2$.
The Lucas residual is zero and the cross and input errors have at least
the retained precision.

For real preparation use the orbit closure of $[h]R_0$. The maps just
displayed are integer multiplication maps of $Z$, including when $t$ or
$m$ is odd. Pulling back the already prepared scalar neighborhood under
$Z\mapsto2Z$ and avoiding the finitely many forbidden point equalities
gives the required compact neighborhoods. Each additional point map is
nonconstant because its integer coefficient is nonzero. The same
finite-index progression recurrence used in the parent preserves these
neighborhoods together with the finite-place conditions. This explains
why introducing a second half-generator multiple does not require a
rational multiplication-by-one-half map on $E$.

The additive contact arguments are still $\lambda,\lambda+T$ with fixed
$T>0$. Rational contact coefficients, coincidences across detectors,
fixed-prime avoidance, and denominator preparation are handled by the
unchanged prime theorem. The factors, trace equations, and both prepared
pole certificates are unchanged. Their local and real limits therefore
remain the parent's limits. The reference and pair charts supply their
rational coordinates, including the new secant slopes. This constructs
a rational witness for every integer $X$. Supports, progressions, root
labels and primes may depend on $(N,X)$; none imposes a further lower
bound on $N$ beyond the fixed scalar threshold.

\section{Literal specification and coordinate accounting}\label{sec:spec}
All indices in this section refer to the parent 204-coordinate unbalanced
table, before the new compaction.

Delete coordinates 169--182 and rows 122--139: the fourth stored root block.
Its only other occurrence is in aggregate row 144. The kept second root
block is rows 104--121. In those rows apply the four rational substitutions
\[
 x_{35}\mapsto(x_{35}-2x_{118})/2,\quad x_{36}\mapsto x_{36}/2,
 \quad x_{37}\mapsto(x_{37}-2x_{118})/2,\quad x_{38}\mapsto x_{38}/2.
\]
These are~\eqref{eq:affine}, not substitutions in the downstream comparison
or Lucas rows.

Before compaction append $(f,h,g,k,s)$ in that order, each in its ordinary
real and $\rho$ component, as coordinates 205--214. Reuse
\[
 F=(x_{60},2x_{61}),\quad G=(x_{29},2x_{30}),\quad
 S=(x_{19},x_{20}),\quad T_H=(x_{183},x_{184}),
\]
\[
 L_4=(x_{43},x_{44}),\qquad L_4^0=(x_{45},x_{46}).
\]
The factors two in the square slots are the parent's half-imaginary
storage convention. Add the nine field quadrics~\eqref{eq:pair} and
\eqref{eq:ratios}, or eighteen rational quadrics.
In aggregate row 144 replace its term
\[
 (x_{173}-x_{175})^2+2(x_{174}-x_{176})^2
\]
by $4\operatorname{Re}(x_4)$, keeping the coefficient multiplying this
term if the primitive row has the opposite sign or a common multiplier.
Everything else in that row and every other retained row is unchanged.

Compact coordinates in increasing order, and put the ten new coordinates
last. The count is
\[
                    204-14+10=200,\qquad 186-18+18=186.
\]
Every resulting row is quadratic. Clear rational denominators in each row,
divide by the gcd of its integer coefficients, and choose the sign so that
the first nonzero coefficient in the specified monomial order is positive.
The defining polynomial is the sum of squares of these 186 rows.
Over $\Q$, this sum vanishes if and only if every row vanishes.

The first component of the new $f$ is coordinate 191 after compaction.
Its fourth power occurs with coefficient one from the real equation
$F=f^2$. No other residual has a pure square in this component: the curve
row uses $fF$, the slope rows use $fg$ or stored squares, the ratio rows
use mixed products, and the aggregate is only linear in $f$. Thus the
coefficient of $x_{191}^4$ in the final polynomial is one for every fixed
parameter. This proves both exact degree four and integer content one.
It does not depend on an expanded monomial count.

The producer and a separately written checker reconstruct the prescribed
changed rows, compare every retained row, and independently expand the
sum of squares. Their receipts, when successfully generated, document
finite algebra only. They do not verify the scalar, modular-height,
prime-pattern, pole-realization or family-height premises, and are not
independent non-author mathematical review.

\section{Balanced coefficients and exact support}\label{sec:balance}
Set $C=149^N$. In the unbalanced rows replace the fixed coefficient
$149^{2N+1}$ by $149C^2$ and replace each of the retained coordinates
$x_{117},x_{118}$ by $C$ times a new coordinate with the same index.
In each row divide by its smallest occurring power of $C$, then take
its primitive integral normalization as in Section~\ref{sec:spec}.
All these transformations are invertible for $C\ne0$. Inspection of the
complete finite row table gives parameter degree at most one and coordinate
degree at most two in every row. Hence no variable or quadratic equation
is added. These balanced rows, rather than the unbalanced rows, may be used
for $P_N$ in Theorem~\ref{thm:main}.

\begin{proposition}[Exact finite coefficient computation]
Put $K_0=31581660^{28}$ and $H_0=(5898050446009651968K_0)^2$.
For every real $C>0$, the balanced sum of squares has 3221 nonzero
coordinate monomials. For each integer $N\ge1$, its coefficient height is
\[
 H(P_N)=\max\{H_0,K_0\,149^{2N}\}.
\]
For $N\ge57$ the second term attains the maximum. This condition concerns
coefficient size only and is not a sufficient defining threshold.
\end{proposition}
\begin{proof}
Expand the 186 primitive quadratic squares and collect equal coordinate
monomials, retaining the parameter as $C$. The complete exact table has
3063 constant profiles, 127 profiles $cC$, 30 profiles $cC^2$, and one
profile $a+bC^2$, with $ab>0$. There are thus 3222 parameter-separated
terms and 3221 coordinate monomials. None vanishes at positive $C$.
The coefficient of $x_{191}^4$ is still one. A pure constant profile has
absolute value $H_0$ (at $x_{112}^4$), and a pure quadratic parameter
profile has absolute value $K_0C^2$ (at $x_{109}^2x_{117}^2$).

The accompanying independent expansion checks, for every pure linear
profile, $c^2\le H_0K_0$. Thus
$|cC|\le\sqrt{H_0K_0C^2}\le\max(H_0,K_0C^2)$.
For every constant/quadratic profile it checks
\[
 |a|K_0+|b|H_0\le H_0K_0,
\]
with absent coefficients taken as zero. Therefore
$|a+bC^2|\le\max(H_0,K_0C^2)$. The two pure profiles show equality.
Finally the exact integer comparison
$149^{114}\ge5898050446009651968^2K_0$ proves the stated branch for
$N\ge57$.

The computations just specified are exhaustive finite integer operations,
not sampling at parameter values. The generator is
\nolinkurl{computation/quartic/shifted200_v33/balance.py}; the separate
\nolinkurl{check_balance.py} replays the coordinate map without importing
the generator, expands every row independently, and checks the displayed
integer inequalities for every profile. The complete tables and their
SHA-256 bindings are distributed with the certificate.
\end{proof}

\begin{thebibliography}{9}
\bibitem{Parent} Project manuscript,
\emph{A quartic in 204 variables from recovered elliptic differences},
\nolinkurl{papers/quartic_existence/quartic204_20261114_v27.tex}.
\bibitem{Weighted} Project weighted-trace manuscript,
\nolinkurl{papers/trace_heights/weighted_trace_20261114_v27.tex}.
\bibitem{Previous} Project manuscript,
\nolinkurl{papers/quartic_existence/quartic212_20261113_v25.tex}.
\bibitem{AllPoints} Project all-point trace-height manuscript,
\nolinkurl{papers/trace_heights/trace_on_all_points_20261113_v25.tex}.
\end{thebibliography}
\end{document}
